% ---
% id: statistics/binomial-box/lem:bernstein-tail
% kind: lem
% title: "Bernstein's inequality for a binomial count"
% summary: >-
%   Either tail of a binomial (n, p) count at distance t > 0 from np has
%   probability at most exp(−t²/(2(v + t/3))) with v = np(1 − p), and every t
%   at least a closed-form root t(L) brings both below e^(−L). The Bernstein
%   endpoint rule evaluates that root.
% prominence: technical
% status: implemented
% uses: [statistics/def:freq]
% assumes: []
% refs: []
% provenance: newly-derived
% implements: certiqs_key.statistics.bernstein:bernstein_deviation
% lean:
% ---

The lemma is Bernstein's inequality for a sum of independent, bounded,
centred random variables, specialised to a binomial count and derived in full
here; its closed-form deviation $t(L)$, the positive root of
$t^{2}-(2L/3)\,t-2vL=0$, is what the Bernstein endpoint rule evaluates.

\begin{lemma}\label{lem:bernstein-tail}
Let $\Ncount{}$ be binomial with parameters $(n,p)$, $n\ge1$, $p\in[0,1]$,
and put $v\coloneqq np(1-p)$. For every real $t>0$,
\begin{equation}\label{eq:bernstein}
  \Pr\big[\Ncount{}-np\ge t\big]\ \le\ \exp\Big(-\frac{t^{2}}{2\,(v+t/3)}\Big),
  \qquad
  \Pr\big[\Ncount{}-np\le -t\big]\ \le\ \exp\Big(-\frac{t^{2}}{2\,(v+t/3)}\Big).
\end{equation}
Consequently, for every real $L>0$ and every
\begin{equation}\label{eq:bernstein-t}
  t\ \ge\ t(L)\ \coloneqq\ \frac{L}{3}+\sqrt{\Big(\frac{L}{3}\Big)^{2}+2vL},
\end{equation}
both tails in \eqref{eq:bernstein} are at most $e^{-L}$.
\end{lemma}

\begin{proof}
\emph{Idea.} The exponential Markov inequality, with the
moment-generating-function bound $\mathbb{E}[e^{\lambda X}]\le
\exp\big(\sigma^{2}(e^{\lambda}-1-\lambda)\big)$ for a centred $X\le1$ and the
estimate $e^{\lambda}-1-\lambda\le\lambda^{2}/(2(1-\lambda/3))$, evaluated at
$\lambda=t/(v+t/3)$; the deviation $t(L)$ then solves a quadratic.

If $p\in\{0,1\}$ then $\Ncount{}=np$ with probability one, both
probabilities in \eqref{eq:bernstein} are $0$ for $t>0$, and there is nothing
to prove. Let $0<p<1$, so $v>0$.

\emph{Step 1 (a moment-generating-function bound).} Let $X$ be a real random
variable with $\mathbb{E}[X]=0$, $\operatorname{Var}(X)=\sigma^{2}<\infty$ and
$X\le1$ almost surely, and let $\lambda>0$. Put
$g(\lambda)\coloneqq e^{\lambda}-1-\lambda$. We claim
\begin{equation}\label{eq:mgf}
  \mathbb{E}\big[e^{\lambda X}\big]\ \le\ 1+\sigma^{2}g(\lambda)\ \le\
  \exp\big(\sigma^{2}g(\lambda)\big).
\end{equation}
The second inequality is $1+z\le e^{z}$. For the first, consider
$\psi(u)\coloneqq(e^{u}-1-u)/u^{2}$ for $u\neq0$, with $\psi(0)\coloneqq1/2$
(its limit). We show that $\psi$ is non-decreasing on $\mathbb{R}$. For
$u\neq0$, $\psi'(u)=N(u)/u^{3}$ with
$N(u)\coloneqq(e^{u}-1)u-2(e^{u}-1-u)=e^{u}(u-2)+u+2$. Now $N(0)=0$,
$N'(u)=e^{u}(u-1)+1$ with $N'(0)=0$, and $N''(u)=u\,e^{u}$, which has the sign
of $u$; hence $N'$ decreases on $(-\infty,0]$ and increases on $[0,\infty)$,
so $N'\ge N'(0)=0$ everywhere, so $N$ is non-decreasing, and with $N(0)=0$
the sign of $N(u)$ is the sign of $u$. Therefore $\psi'(u)=N(u)/u^{3}\ge0$ for
every $u\neq0$, and $\psi$ is non-decreasing (continuity at $0$ follows from
the series $\psi(u)=\sum_{k\ge2}u^{k-2}/k!$). For a real $x\le1$ and
$\lambda>0$ we have $\lambda x\le\lambda$, hence
$\psi(\lambda x)\le\psi(\lambda)$, i.e.\ for $x\neq0$
\begin{equation*}
  \frac{e^{\lambda x}-1-\lambda x}{\lambda^{2}x^{2}}\ \le\
  \frac{e^{\lambda}-1-\lambda}{\lambda^{2}},
  \qquad\text{so}\qquad
  e^{\lambda x}\ \le\ 1+\lambda x+x^{2}g(\lambda),
\end{equation*}
and the last inequality holds at $x=0$ too. Apply it to $X$ pointwise and
take expectations: $\mathbb{E}[e^{\lambda X}]\le1+\lambda\,\mathbb{E}[X]
+g(\lambda)\,\mathbb{E}[X^{2}]=1+\sigma^{2}g(\lambda)$, which is
\eqref{eq:mgf}.

\emph{Step 2 (exponential Markov).} Write $\Ncount{}=\sum_{j=1}^{n}\chi_j$ with
$\chi_j$ independent Bernoulli$(p)$ (Definition~\ref{def:freq}), and
$X_j\coloneqq\chi_j-p$. Then $\mathbb{E}[X_j]=0$,
$\operatorname{Var}(X_j)=p(1-p)$, and $X_j\le1-p\le1$. For $\lambda>0$,
Markov's inequality applied to the non-negative variable
$e^{\lambda\sum_jX_j}$ and independence give
\begin{equation*}
  \Pr\Big[\sum_{j}X_j\ge t\Big]
  \ \le\ e^{-\lambda t}\,\mathbb{E}\Big[e^{\lambda\sum_jX_j}\Big]
  \ =\ e^{-\lambda t}\prod_{j=1}^{n}\mathbb{E}\big[e^{\lambda X_j}\big]
  \ \le\ \exp\big(-\lambda t+v\,g(\lambda)\big),
\end{equation*}
by \eqref{eq:mgf} with $\sigma^{2}=p(1-p)$ and $np(1-p)=v$.

\emph{Step 3 (an elementary bound on $g$).} For integers $k\ge2$,
$k!\ge2\cdot3^{k-2}$ (induction: $2!=2$, and
$(k+1)!=(k+1)\,k!\ge3\cdot2\cdot3^{k-2}$). Hence for $0<\lambda<3$,
\begin{equation*}
  g(\lambda)=\sum_{k\ge2}\frac{\lambda^{k}}{k!}
  \ \le\ \frac{\lambda^{2}}{2}\sum_{k\ge2}\Big(\frac{\lambda}{3}\Big)^{k-2}
  \ =\ \frac{\lambda^{2}}{2\,(1-\lambda/3)}.
\end{equation*}

\emph{Step 4 (the choice of $\lambda$).} Put
$\lambda\coloneqq t/(v+t/3)$, which lies in $(0,3)$ because $v>0$. Then
$1-\lambda/3=v/(v+t/3)$, so Step 3 gives
$g(\lambda)\le\lambda^{2}(v+t/3)/(2v)=t^{2}/\big(2v(v+t/3)\big)$, and the
exponent of Step 2 is at most
\begin{equation*}
  -\lambda t+v\,g(\lambda)
  \ \le\ -\frac{t^{2}}{v+t/3}+\frac{t^{2}}{2(v+t/3)}
  \ =\ -\frac{t^{2}}{2(v+t/3)},
\end{equation*}
which is the first inequality of \eqref{eq:bernstein}. For the second, apply
Steps 1--4 to $-X_j=p-\chi_j$, which has the same mean and variance and
satisfies $-X_j\le p\le1$: $\Pr[\sum_j(-X_j)\ge t]=\Pr[\Ncount{}-np\le-t]$.

\emph{Step 5 (the closed-form deviation).} For $t>0$, the inequality
$t^{2}/(2(v+t/3))\ge L$ is equivalent to $t^{2}-(2L/3)\,t-2vL\ge0$. The
quadratic on the left has the two real roots $L/3\pm\sqrt{(L/3)^{2}+2vL}$, of
which $t(L)$ is the positive one and the other is $\le0$; so the inequality
holds for every $t\ge t(L)$, and then both bounds in \eqref{eq:bernstein}
are at most $e^{-L}$.
\end{proof}
